% Template TCC abnTeX2 — pesquisador-br-skill
% Compatível com classe abntex2 (recomendada para TCC/dissertação/tese ABNT)
% Compilação: pdflatex + bibtex + pdflatex + pdflatex
% Documentação: https://www.abntex.net.br

\documentclass[
    12pt,                       % corpo 12pt
    openright,                  % capítulos abrem em página direita
    twoside,                    % impressão frente e verso
    a4paper,                    % papel A4
    chapter=TITLE,              % capítulos em maiúsculas
    sumario=tradicional,        % sumário NBR 6027
    section=TITLE,
    english,                    % inglês para abstract
    spanish,                    % espanhol opcional
    brazil                      % última = idioma principal
]{abntex2}

% ─── Pacotes ────────────────────────────────────────
\usepackage{cmap}                           % suporte a copy-paste
\usepackage[utf8]{inputenc}                 % UTF-8
\usepackage[T1]{fontenc}                    % fontes europeias
\usepackage{lastpage}                       % páginas totais
\usepackage{indentfirst}                    % indenta primeiro parágrafo
\usepackage{color}                          % cores
\usepackage{graphicx}                       % imagens
\usepackage{microtype}                      % melhorias tipográficas
\usepackage[brazilian,hyperpageref]{backref}
\usepackage[alf]{abntex2cite}               % citações ABNT autor-data

% ─── Configurações ──────────────────────────────────
\titulo{TÍTULO DO TRABALHO: subtítulo se houver}
\autor{Nome Completo Sobrenome}
\local{Manaus -- AM}
\data{2026}
\orientador{Prof.\textsuperscript{a} Dr.\textsuperscript{a} Nome do Orientador}
% \coorientador{Prof. Dr. Nome do Coorientador}
\instituicao{Universidade Federal do Amazonas \par
            Instituto de [...] \par
            Curso de [...]}
\tipotrabalho{Trabalho de Conclusão de Curso}
\preambulo{Trabalho de Conclusão de Curso apresentado ao Curso de
[...] da Universidade Federal do Amazonas, como requisito parcial
para a obtenção do título de Bacharel/Licenciado em [...].}

% ─── Hyperref ───────────────────────────────────────
\usepackage{hyperref}
\hypersetup{
    pdftitle={\@title},
    pdfauthor={\@author},
    pdfsubject={\imprimirpreambulo},
    pdfkeywords={tcc}{abntex2}{universidade},
    colorlinks=true,
    linkcolor=blue,
    citecolor=blue,
    filecolor=magenta,
    urlcolor=blue,
}

% ═══════════════════════════════════════════════════
\begin{document}
\selectlanguage{brazil}
\frenchspacing

% ─── PRÉ-TEXTUAIS ──────────────────────────────────
\pretextual

\imprimircapa
\imprimirfolhaderosto

% Errata (opcional)
% \begin{errata}
% Errata aqui
% \end{errata}

% Folha de aprovação (manual)
\begin{folhadeaprovacao}
\begin{center}
{\ABNTEXchapterfont\large\imprimirautor}

\vspace*{\fill}\vspace*{\fill}
\begin{center}
\ABNTEXchapterfont\bfseries\Large\imprimirtitulo
\end{center}
\vspace*{\fill}

\hspace{.45\textwidth}\begin{minipage}{.5\textwidth}
\imprimirpreambulo
\end{minipage}\vspace*{\fill}
\end{center}

Trabalho aprovado. Manaus -- AM, \today:

\assinatura{\textbf{\imprimirorientadorRotulo} \\ \imprimirorientador}
\assinatura{\textbf{Prof. Dr. Nome 1} \\ Instituição}
\assinatura{\textbf{Prof. Dr. Nome 2} \\ Instituição}
\end{folhadeaprovacao}

% Dedicatória (opcional)
\begin{dedicatoria}
\vspace*{\fill}
\centering
\noindent
\textit{Dedicatória ao [pessoa/grupo].}
\vspace*{\fill}
\end{dedicatoria}

% Agradecimentos (opcional)
\begin{agradecimentos}
Agradeço à orientadora, à banca, à CAPES (processo nº \textit{X}),
ao CNPq (processo nº \textit{Y}), e à minha família.
\end{agradecimentos}

% Epígrafe (opcional)
\begin{epigrafe}
\vspace*{\fill}
\begin{flushright}
\textit{``Citação que inspira o trabalho.''} \\
(AUTOR, 2020, p. 25)
\end{flushright}
\end{epigrafe}

% Resumo
\setlength{\absparsep}{18pt}
\begin{resumo}
\noindent O presente trabalho objetivou [objetivo geral]. Adotou-se
[tipo de pesquisa]. Foram coletados [dados]. Os resultados indicaram
[achado]. Conclui-se que [síntese].

\noindent\textbf{Palavras-chave}: termo 1; termo 2; termo 3;
termo 4; termo 5.
\end{resumo}

% Abstract
\begin{resumo}[Abstract]
\begin{otherlanguage*}{english}
\noindent This work objective was [...]. The methodology adopted
was [...]. Data were collected through [...]. Results indicated
[...]. It is concluded that [...].

\noindent\textbf{Keywords}: term 1; term 2; term 3; term 4; term 5.
\end{otherlanguage*}
\end{resumo}

% Listas (figuras, tabelas, abreviaturas)
\pdfbookmark[0]{\listfigurename}{lof}
\listoffigures*
\cleardoublepage

\pdfbookmark[0]{\listtablename}{lot}
\listoftables*
\cleardoublepage

% Sumário
\pdfbookmark[0]{\contentsname}{toc}
\tableofcontents*
\cleardoublepage

% ─── TEXTUAIS ──────────────────────────────────────
\textual

% Capítulo 1: Introdução
\chapter{INTRODUÇÃO}
\label{cap:intro}

\section{Contexto e motivação}
[Texto da contextualização. Citação como exemplo: \citeonline{gil2017}.]

\section{Problema de pesquisa}
[Pergunta de pesquisa.]

\section{Objetivos}
\subsection{Objetivo geral}
[Verbo no infinitivo + objeto + condição].

\subsection{Objetivos específicos}
\begin{enumerate}[label=\alph*)]
\item [Verbo + objeto];
\item [Verbo + objeto];
\item [Verbo + objeto].
\end{enumerate}

\section{Justificativa}
[Texto.]

\section{Estrutura do trabalho}
[Texto.]

% Capítulo 2: Referencial Teórico
\chapter{REFERENCIAL TEÓRICO}
\label{cap:teorico}

\section{[Conceito-chave 1]}
[Discussão. Exemplo de citação: \cite{minayo2014}.]

\section{[Conceito-chave 2]}
[Discussão.]

% Capítulo 3: Metodologia
\chapter{METODOLOGIA}
\label{cap:metodo}

\section{Tipo e abordagem}
A pesquisa caracteriza-se como [tipo]\dots
\citeonline{gil2017}.

\section{Sujeitos / corpus}
[Descrição.]

\section{Procedimentos de coleta}
[Descrição.]

\section{Procedimentos de análise}
[Descrição. Análise de Conteúdo conforme \citeonline{bardin2011}.]

\section{Aspectos éticos}
[CAAE, TCLE, LGPD.]

% Capítulo 4: Resultados e Discussão
\chapter{RESULTADOS E DISCUSSÃO}
\label{cap:resultados}

\section{[Eixo temático 1]}
[Apresentação dos resultados.]

\section{[Eixo temático 2]}
[Apresentação.]

% Capítulo 5: Considerações finais
\chapter{CONSIDERAÇÕES FINAIS}
\label{cap:conclusao}

[Síntese, contribuições, limitações, agenda futura.]

% ─── PÓS-TEXTUAIS ──────────────────────────────────
\postextual

% Referências
\bibliography{referencias}

% Apêndices
\begin{apendicesenv}
\partapendices

\chapter{ROTEIRO DE ENTREVISTA}
[Conteúdo do apêndice.]

\end{apendicesenv}

% Anexos
\begin{anexosenv}
\partanexos

\chapter{PARECER COMITÊ DE ÉTICA}
[Conteúdo do anexo.]

\end{anexosenv}

\end{document}
